\section{Local theory and continuation for finite-energy fields}
\label{sec:continuation-proof}

We prove Proposition~\ref{prop:continuation}.  The distinction between
bounded smooth derivatives and square-integrable derivatives matters here:
the initial fields in the theorem need not belong to $H^5$.
We first give the local theory for Sobolev fields, then replace the fields
outside a sufficiently large ball while preserving the divergence
constraints.  For any prescribed finite horizon, a vacuum field
transfers solutions between the original and modified data on every
interval of existence within that horizon.

Throughout this section, $F=(E,B)$,
$u(v)=v/(1+|v|^2)^{1/2}$, and
\[
 \rho_f=\int_{\R^3}f\,dv,
 \qquad j_f=\int_{\R^3}u(v)f\,dv.
\]
The data are $0\le f_0\in C_c^\infty(\R^6)$ and
$E_0,B_0\in C_b^\infty(\R^3)\cap L^2(\R^3)$, with every spatial
derivative bounded and with
$\nabla\cdot E_0=\rho_{f_0}$, $\nabla\cdot B_0=0$.
Choose $R_0,R_v$ so that
\[
 \operatorname{supp}f_0\subset
 \{(x,v):|x|\le R_0,\ |v|\le R_v\}.
\]
As in Proposition~\ref{prop:continuation}, a classical solution is a
$C^1$ pointwise solution whose fields are continuous into $L^2_x$ and whose
density has compact phase-space support on every compact time interval.

\subsection{Characteristics and the Maxwell evolution}

For any $C^1$ fields on a closed finite time interval, the characteristic
system is
\begin{equation}\label{cont:characteristics}
 \dot X=u(V),\qquad
 \dot V=E(t,X)+u(V)\times B(t,X).
\end{equation}
It has a unique solution through every phase point throughout that
interval.  Indeed, $|\dot X|\le1$ confines $X$ to a compact spatial ball;
the fields are bounded on the resulting compact spacetime cylinder, so
$|\dot V|$ is bounded there.  The usual local ODE solution therefore
cannot escape in finite time.  Transport along these characteristics gives
\begin{equation}\label{cont:spatial-support}
 f(t,X(t),V(t))=f_0(X(0),V(0)),\qquad
 \operatorname{supp}_x f(t)\subset\{|x|\le R_0+t\}.
\end{equation}
On every closed finite slab the same argument gives a common momentum
bound, provided the $C^1$ solution exists on that closed slab.  It does
not give a bound up to an excluded endpoint of existence.
The phase-space vector field in \eqref{cont:characteristics} has zero
divergence: $u=\nabla_v(1+|v|^2)^{1/2}$ and hence
$\nabla_v\cdot(u\times B)=0$.
Consequently positivity and the $L^\infty$ norm of $f$ are preserved.
Integration of the transport equation gives
$\partial_t\rho_f+\nabla\cdot j_f=0$, which, together with Maxwell's
equations, propagates both divergence constraints.

We shall use the homogeneous \emph{first-order Maxwell evolution} on
$F=(E,B)$.  If $C_\xi w=\xi\times w$, its Fourier generator is
\[
 A(\xi)=\begin{pmatrix}0&iC_\xi\\-iC_\xi&0\end{pmatrix},
 \qquad A(\xi)^*=-A(\xi).
\]
Thus $U(t)$, defined by the multiplier $\exp(tA(\xi))$, is a strongly
continuous unitary group on $L^2$ and on every $H^k$.
No divergence constraint is needed for this assertion.  The formula
\begin{equation}\label{cont:maxwell-duhamel}
 F(t)=U(t)F(0)+\int_0^t U(t-s)(-j_f(s),0)\,ds
\end{equation}
therefore yields
\[
 \|F(t)\|_{H^k}\le\|F(0)\|_{H^k}
       +\int_0^t\|j_f(s)\|_{H^k}\,ds.
\]

We also need finite propagation without any global integrability
assumption.  For a $C^1$ Maxwell field with current $j$, set
$e=(|E|^2+|B|^2)/2$.  Direct differentiation gives
\begin{equation}\label{cont:local-energy}
 \partial_t e+\nabla\cdot(E\times B)=-E\cdot j,
 \qquad |E\times B|\le e.
\end{equation}
Integrating over the shrinking balls $B(x,r+t-s)$, $0\le s\le t$,
shows that a field with zero initial data and zero current in this
backward cone vanishes on the final ball $B(x,r)$: the boundary term is
$-\int_{\partial B}(e+(E\times B)\cdot\nu)\le0$.
This local argument uses neither global $L^2$ bounds nor divergence
constraints.

\begin{lemma}[Exterior comparison]\label{cont:exterior}
Suppose that a $C^1$ solution exists on $[0,S]$, that $f_0$ is compactly
supported as above, and that $F(0)\in H^k(\R^3)$ for every integer $k$.
Then the fields belong to $C([0,S];L^2)$ and have bounded first spatial
derivatives on the entire slab.  More precisely,
\[
 F(t,x)=U(t)F(0)(x)\qquad\text{if }|x|>R_0+t.
\]
\end{lemma}
\begin{proof}
The characteristic argument gives compact phase-space support on the
whole slab, so $j_f$ is $C^1$ and vanishes outside $|x|\le R_0+t$.
The homogeneous field $F^{\rm h}(t)=U(t)F(0)$ is smooth and has bounded
spatial derivatives on $[0,S]\times\R^3$, by its uniform Sobolev bounds.
The difference $W=F-F^{\rm h}$ has zero initial data and current $j_f$.
If $|x|>R_0+t$, choose $r>0$ smaller than $|x|-R_0-t$.
Every point $(s,y)$ of the backward cone from $B(x,r)$ satisfies
$|y|>R_0+s$, so its current is zero.  The local energy argument gives
$W(t)=0$ on $B(x,r)$.
Thus $W$ is supported in the fixed ball $|x|\le R_0+S$ on this slab.
Its $C^1$ regularity on that compact cylinder gives
$W\in C([0,S];L^2)$ and uniform bounds on its first derivatives.
Adding $F^{\rm h}$ proves the claims.
Notice that $F^{\rm h}$ can have nonzero electric divergence; only the
first-order Maxwell equations, not a component free-wave equation, were
used for this comparison.
\end{proof}

\subsection{Classical uniqueness and local Sobolev solutions}
\label{cont:local}

Uniqueness holds in the full classical class stated above.  To see this,
let $(f_i,F_i)$, $i=1,2$, have the same data, and work on an arbitrary
compact common time interval.  Choose $R$ containing both momentum
supports, and put $g=f_1-f_2$, $\delta F=F_1-F_2$.
With the transport operator of the first solution,
\[
 (\partial_t+u\cdot\nabla_x+(E_1+u\times B_1)\cdot\nabla_v)g
 =-(\delta E+u\times\delta B)\cdot\nabla_v f_2.
\]
The phase-space divergence is zero, $g$ is compactly supported, and
\[
 \|j_{f_1}-j_{f_2}\|_{L^2_x}\le C_R\|g\|_{L^2_{x,v}}.
\]
Since $\nabla_v f_2$ is bounded on this slab, integration of the transport
and Maxwell energies gives, with a slab-dependent finite constant $C$,
\begin{equation}\label{cont:uniqueness-energy}
 \|g(t)\|_2^2+\|\delta F(t)\|_2^2
 \le C\int_0^t\bigl(\|g(s)\|_2^2+\|\delta F(s)\|_2^2\bigr)\,ds.
\end{equation}
For completeness, the Maxwell integration is justified with cutoffs
$\chi(x/N)$, where $\chi=1$ on the unit ball and
$|\nabla\chi(x/N)|\le C/N$.  Its boundary error is at most
$C N^{-1}\int_0^S\|\delta F(s)\|_2^2\,ds$, which tends to zero because
$\delta F\in C_tL^2_x$.  No square-integrable field derivatives are
required.  Gronwall's inequality in \eqref{cont:uniqueness-energy}
proves uniqueness.

We next record the estimate needed both to construct smooth solutions
and to preserve their regularity.  For an integer $k\ge6$, write
\[
 Y_k=\|f\|_{H^k(\R^6)}+\|F\|_{H^k(\R^3)}.
\]
For a smooth solution with momentum support in $|v|\le R$,
\begin{equation}\label{cont:tame}
 Y_k(t)\le Y_k(0)+C_{k,R}\int_0^t
 \bigl(1+\|F(s)\|_{W^{1,\infty}_x}
          +\|\nabla_{x,v}f(s)\|_\infty\bigr)Y_k(s)\,ds.
\end{equation}
The $H^k$ norm of $F$ is the Hilbert norm of the six-component vector
$(E,B)$; this is the norm preserved by $U(t)$.
The coefficient in \eqref{cont:tame} involves no derivatives above
first order.  For smooth solutions in these Sobolev spaces, bounds on
the momentum support and this coefficient give a bound on $Y_k$ by
Gronwall's inequality.  The local construction below will use the
$Y_6$ bound to obtain a common restart time.

To retain this dependence in the commutators, we use the integer
Gagliardo--Nirenberg inequality on the whole space and include its
proof; see also
\cite[Lecture~II]{Nirenberg1959} for its classical context.
For $d\ge1$, an integer $s\ge1$, a scalar or finite-dimensional
vector-valued $g\in H^s(\mathbb R^d)\cap L^\infty$, and an integer $0<j<s$,
\begin{equation}\label{cont:integer-interpolation}
 \|D^j g\|_{L^{2s/j}}
 \le C_{d,s}\|g\|_\infty^{1-j/s}
                   \|D^s g\|_2^{j/s}.
\end{equation}
Here $D^j g$ denotes the array of all ordered $j$th derivatives, with
its Euclidean pointwise norm; this convention is equivalent to the
usual multiindex convention up to constants depending on $d,s$.
For $g\in C_c^\infty$ put $p_j=2s/j$ and
$A_j=\|D^jg\|_{p_j}$, with $A_0=\|g\|_\infty$.
Integrating one derivative by parts in
$\int|D^jg|^{p_j}$ gives
\[
 A_j^{p_j}
 \le C_{d,s}\int |D^{j-1}g|\,|D^{j+1}g|\,
                         |D^jg|^{p_j-2}
 \le C_{d,s} A_{j-1}A_{j+1} A_j^{p_j-2}.
\]
The last step is H\"older's inequality, since
$1/p_{j-1}+1/p_{j+1}+(p_j-2)/p_j=1$, with $p_0=\infty$.
Thus $A_j^2\le C_{d,s}A_{j-1}A_{j+1}$.
For a nonzero compactly supported $g$ all $A_j$ are positive.
Choose $C\ge1$ valid for these inequalities and define
\[
 b_j=\log A_j-(1-j/s)\log A_0-(j/s)\log A_s.
\]
Then $b_0=b_s=0$ and
$2b_j-b_{j-1}-b_{j+1}\le\log C$.
The discrete maximum principle, applied after subtracting
$\tfrac12(\log C)j(s-j)$, gives
$b_j\le\tfrac12(\log C)j(s-j)$, which proves
\eqref{cont:integer-interpolation}; the zero function is immediate.
To pass to $H^s\cap L^\infty$, first multiply by smooth compactly
supported cutoffs bounded by one and then mollify with nonnegative
unit-mass kernels. A diagonal sequence converges in $H^s$, has
$L^\infty$ norm at most $\|g\|_\infty$, and has the relevant derivatives
converging almost everywhere along a subsequence. Fatou's lemma proves
\eqref{cont:integer-interpolation} in the stated class.

We give the commutator argument explicitly.  If $s\ge1$ is an integer,
$|\alpha|+|\beta|=s$, and the functions are in $H^s\cap L^\infty$,
then the Gagliardo--Nirenberg interpolation inequality and H\"older's
inequality give
\begin{equation}\label{cont:tame-product}
 \|D^\alpha g\,D^\beta h\|_2
 \le C_s\bigl(\|g\|_\infty\|h\|_{H^s}
                  +\|h\|_\infty\|g\|_{H^s}\bigr).
\end{equation}
Indeed, put $j=|\alpha|$. For $0<j<s$ apply
\eqref{cont:integer-interpolation} to $g$
and the analogous bound with derivative order $s-j$ to $h$, then use
weighted arithmetic--geometric mean.  The cases $j=0,s$ are direct.
The same estimate for lower positive total derivative orders follows
by using that order in place of $s$; total order zero is immediate.

Differentiate the Vlasov equation in $(x,v)$ to order at most $k$ and
pair with the corresponding derivative of $f$ in $L^2$.
The leading transport term vanishes by phase-space divergence.
Derivatives falling on $u$ are bounded and cost only $\|f\|_{H^k}$.
For the force terms, choose a smooth momentum cutoff $\zeta_R$ equal
to one on a neighborhood of $|v|\le R$, and set
$a_R=\zeta_R(v)(E+u\times B)$.
It is used only in the commutators; the leading transport operator is
unchanged.  We have
\[
 \|a_R\|_{H^k(\R^6)}\le C_{k,R}\|F\|_{H^k(\R^3)},\qquad
 \|\nabla a_R\|_\infty\le C_R\|F\|_{W^{1,\infty}_x}.
\]
Each commutator product is a derivative of $\nabla a_R$ times a
derivative of $\nabla f$, with total derivative order at most $k-1$.
Equation~\eqref{cont:tame-product} therefore bounds their sum by
\[
 C_{k,R}\bigl(\|F\|_{W^{1,\infty}_x}\|f\|_{H^k}
             +\|F\|_{H^k}\|\nabla f\|_\infty\bigr).
\]
Finally, $\|j_f\|_{H^k_x}\le C_R\|f\|_{H^k_{x,v}}$ by
Cauchy--Schwarz in $v$, and \eqref{cont:maxwell-duhamel} completes
\eqref{cont:tame}.  The estimates for unsquared norms follow from the
squared energy estimates by regularizing a vanishing norm and taking
the regularization parameter to zero.

Here are sufficient details of the local construction to make the
restart time in this argument explicit.  Suppose first that $F_0$
belongs to every $H^k$.  Starting with the time-independent data, solve
successively the linear transport equation for $f^{n+1}$ with force
$F^n$, and the linear Maxwell equations for $F^{n+1}$ with current
$j_{f^n}$; every iterate has the same initial data.
Characteristics and \eqref{cont:maxwell-duhamel} produce smooth iterates.
Because $H^6(\R^6)$ controls first derivatives in $L^\infty$, the
preceding estimates give, while all momentum supports lie in a fixed
ball,
\[
 Y_6^{n+1}(t)\le Y_6(0)+C_R\int_0^t
       \bigl((1+Y_6^n)Y_6^{n+1}+Y_6^n\bigr)\,ds.
\]
Choose $M>2(Y_6(0)+1)$, and then choose $\tau>0$ small in terms of
$M$ and $R_v$.  Induction gives $Y_6^n\le M$ on $[0,\tau]$;
shrinking $\tau$ further makes the momentum increase at most one,
since $|\dot V|\le |E^n|+|B^n|\le C M$.
This closes the support condition with $R=R_v+1$.

The $L^2$ difference estimate for consecutive iterates has the form
\[
 D_{n+1}(t)\le C_{R,M}\int_0^t D_n(s)\,ds,
 \qquad
 D_n=\|f^n-f^{n-1}\|_2+\|F^n-F^{n-1}\|_2.
\]
The term involving $f^n$ uses its uniform first-derivative bound,
exactly as in \eqref{cont:uniqueness-energy}; the Maxwell estimate is
linear.  Taking $C_{R,M}\tau<1/2$ makes the iterates Cauchy in
$C([0,\tau];L^2)$.
For every higher $k$, the uniform $C^1$ bounds already obtained give
\[
 Y_k^{n+1}(t)\le Y_k(0)+C_{k,R,M}\int_0^t
                    (Y_k^{n+1}(s)+Y_k^n(s))\,ds.
\]
The exponential $Y_k(0)e^{2C_{k,R,M}t}$ is a common upper bound, by
induction and the scalar comparison inequality.  Interpolating the
$L^2$ convergence with the uniform $H^{k+1}$ bounds gives convergence
in $C([0,\tau];H^k)$ for every $k$.  The limit solves the equations;
it preserves positivity, compact phase-space support, and the
constraints.  Sobolev embedding and the equations give smoothness in
both space and time.  Crucially, the common lifespan $\tau$ depends
only on $Y_6(0)$ and the initial momentum radius, not on higher norms.

\subsection{Changing the fields outside the particle region}

We first construct the vacuum evolution needed for data whose bounded
derivatives need not lie in $L^2$.

\begin{lemma}[Smooth finite-energy vacuum fields]\label{cont:vacuum}
Let $D_E,D_B\in C_b^\infty(\R^3)\cap L^2(\R^3)$ be divergence-free.
There is a global smooth vacuum Maxwell solution on $[0,\infty)$ with these data,
continuous into $L^2$, whose spatial derivatives of every order are
bounded on each finite slab.  If the data vanish on $B(0,L)$, the
solution vanishes wherever $|x|+t<L$.
\end{lemma}
\begin{proof}
For $M_tg(x)=(4\pi)^{-1}\int_{S^2}g(x+t\omega)\,d\omega$, define
\begin{align*}
 E^{\rm v}(t)&=\partial_t(tM_tD_E)+tM_t(\nabla\times D_B),\\
 B^{\rm v}(t)&=\partial_t(tM_tD_B)-tM_t(\nabla\times D_E).
\end{align*}
Kirchhoff's formula gives smooth free-wave solutions with all the
asserted finite-slab derivative bounds.  Their Maxwell residuals
$E^{\rm v}_t-\nabla\times B^{\rm v}$ and
$B^{\rm v}_t+\nabla\times E^{\rm v}$ solve the wave equation with
zero initial values.  Their initial time derivatives are respectively
$\nabla(\nabla\cdot D_E)$ and $\nabla(\nabla\cdot D_B)$, hence zero.
Local wave uniqueness makes both residuals vanish.  The divergences
vanish as well.  The spherical-mean formula proves the stated finite
propagation.

To justify strong $L^2$ continuity without assuming square-integrable
derivatives, convolve the data with a smooth compactly supported
approximate identity $\eta_\epsilon$.  The resulting divergence-free
data belong to every $H^k$, converge in $L^2$, and converge uniformly
with every fixed spatial derivative: the next bounded derivative makes
each derivative uniformly continuous.  Their Kirchhoff solutions
converge with all fixed derivatives on finite slabs to the displayed
solution.  The same solutions are $U(t)(\eta_\epsilon*D_E,
\eta_\epsilon*D_B)$, whose $L^2$ convergence is uniform in $t$ by
unitarity.  The two limits agree locally, so the displayed solution is
$U(t)(D_E,D_B)$.  It is therefore continuous into $L^2$ and preserves
its $L^2$ norm.
\end{proof}

Fix a finite horizon $H>0$.  Set $\rho_0=\rho_{f_0}$ and take its
Coulomb field
\begin{equation}\label{cont:coulomb}
 E_{\rm C}=-\nabla(-\Delta)^{-1}\rho_0,
 \qquad \nabla\cdot E_{\rm C}=\rho_0.
\end{equation}
This field belongs to every $H^s$, $s\ge0$.  Indeed, its Fourier
transform is $-i\xi|\xi|^{-2}\widehat\rho_0(\xi)$.
Near zero its size is $O(|\xi|^{-1})$, whose square is integrable in
three dimensions, and at infinity $\widehat\rho_0$ decays faster than
any power.  Sobolev embedding gives bounded smooth derivatives.
This argument allows arbitrary total charge.

For a smooth divergence-free vector field $G$ on $\R^3$, define
\begin{equation}\label{cont:potential}
 \mathcal A_G(x)=\left(\int_0^1 tG(tx)\,dt\right)\times x.
\end{equation}
Then $\nabla\times\mathcal A_G=G$.  To verify this, put
$V(x)=\int_0^1tG(tx)\,dt$.  We have $\nabla\cdot V=0$ and
\[
 \nabla\times(V\times x)=2V+(x\cdot\nabla)V
 =\int_0^1\frac{d}{dt}\bigl(t^2G(tx)\bigr)\,dt=G(x).
\]
Choose $\chi\in C_c^\infty(\R^3)$ equal to one on
$B(0,L)$, where $L=R_0+2H+1$, and set
\begin{equation}\label{cont:modified-data}
 \begin{split}
 \overline E_0&=E_{\rm C}+
     \nabla\times\bigl(\chi\mathcal A_{E_0-E_{\rm C}}\bigr),\\
 \overline B_0&=\nabla\times\bigl(\chi\mathcal A_{B_0}\bigr).
 \end{split}
\end{equation}
These fields belong to every $H^k$, satisfy the original divergence
constraints, and agree with $E_0,B_0$ on $B(0,L)$.
Their difference from the original fields is divergence-free,
smooth with bounded derivatives, and in $L^2$.
Let $D(t)$ be the vacuum solution of Lemma~\ref{cont:vacuum} with
initial data
$D(0)=(E_0-\overline E_0,B_0-\overline B_0)$.
Then
\begin{equation}\label{cont:vacuum-vanishing}
 D(t,x)=0\quad\text{if}\quad
 0\le t\le H,\quad |x|\le R_0+t,
\end{equation}
because $|x|+t\le R_0+2H<L$.
Consequently
\begin{equation}\label{cont:field-equivalence}
 (f,F)\longleftrightarrow(f,F-D)
\end{equation}
is an equivalence, up to time $H$, between classical solutions with
the original and modified initial fields.
Maxwell's equations and both constraints are preserved by subtracting
a divergence-free vacuum solution.  The transport equation is
unchanged on the particle region by \eqref{cont:vacuum-vanishing};
outside that region $f$ and its momentum derivatives vanish.
The equivalence preserves $C_tL^2$ and compact phase-space support.

The Sobolev construction above, applied to \eqref{cont:modified-data}
with, for example, $H=1$, and followed by addition of $D$, now proves
local classical existence for the full initial field class.
The uniqueness already proved permits these local solutions to be
joined into a unique maximal classical solution.

\subsection{Smoothness through a classical interval}
\label{cont:persistence}

Suppose first that the initial fields belong to every $H^k$ and that
a classical solution exists on $[0,S]$.
Lemma~\ref{cont:exterior} gives a uniform global
$W^{1,\infty}_x$ bound for its fields on this slab.
Compact phase-space support and $C^1$ regularity give uniform bounds
for the momentum radius and $\nabla_{x,v}f$.
Wherever the solution is smooth, \eqref{cont:tame} therefore bounds
every $Y_k$, in particular $Y_6$, by a constant depending only on this
slab and the initial norm.

To continue smoothness, suppose a first loss could occur before or at
$S$.  Choose times increasing to that first time from below.
The $Y_6$ bounds and momentum bounds just obtained give the same
positive local lifespan for the Sobolev construction at all these
times.  Restarting sufficiently close to the alleged loss extends a
smooth solution beyond it.  Classical uniqueness identifies this
solution with the given one on their common interval, a contradiction.
This argument requires no prior existence of an $H^6$ limit at the
endpoint.  The higher estimates on the same lifespan give persistence
of every Sobolev order, and the equations give all time derivatives.

For the original field class, fix any compact classical interval,
choose a larger finite horizon in \eqref{cont:modified-data}, and
subtract the associated vacuum field.  The modified solution has
all-Sobolev initial fields, so the preceding argument applies.
Adding back the smooth vacuum proves smoothness on the chosen slab
for the original solution.

The solutions just constructed satisfy the hypotheses of
Lemma~\ref{ret:energy}, so their total energy is conserved.
Here $C_tL^2$ was established independently; no square-integrable
field derivatives are needed to apply that energy identity.

\subsection{The bounded-momentum criterion}
\label{cont:continuation}

The continuation input is the Glassey--Strauss criterion in
\cite[Theorem~1.1 and Footnote~1, arXiv version]{LukStrain2014Continuation}.
In its normalization, nonnegative compactly supported $H^5$ particle
data and $H^5$ field data satisfying the constraints have the following
property: a unique classical solution on $[0,T)$, $T<\infty$, extends
uniquely as a $C^1$ solution past $T$ if there is a bounded continuous
function $\kappa:[0,T)\to(0,\infty)$ such that $f(t,x,v)=0$ whenever
$|v|\ge\kappa(t)$.
The $H^5$ formulation in that statement does not require a separate
momentum bound for approximation iterates.
Its density and current are $4\pi\int f_{\rm LS}\,dv$ and
$4\pi\int u f_{\rm LS}\,dv$, so it applies to the present normalization
with
\begin{equation}\label{cont:normalization}
 f_{\rm LS}=f/(4\pi),\qquad E_{\rm LS}=E,
 \qquad B_{\rm LS}=B.
\end{equation}

Now let a classical solution for the original data exist on $[0,T)$,
where $T<\infty$, and suppose its momentum support is uniformly
bounded there by $P_*<\infty$.
Perform \eqref{cont:modified-data} with $H=T+1$.
The modified solution $(f,F-D)$ has all-Sobolev initial fields;
its density is smooth and compactly supported initially, and the
constraints and uniqueness have already been checked.
After \eqref{cont:normalization}, all hypotheses of the cited theorem
hold with the constant continuous choice $\kappa=P_*+1$.
It gives a $C^1$ extension of the modified solution to $[0,T+\epsilon)$;
shrink $\epsilon$ so that $\epsilon<1$.

It remains to check that this extension is in the classical class used
here.  On any closed subinterval of $[0,T+\epsilon)$, its
characteristics give compact phase-space support, as in
\eqref{cont:spatial-support}.  Lemma~\ref{cont:exterior} supplies
$C_tL^2$ for its fields, including continuity through $T$.
Thus it is a classical solution in the stated class.
Adding back $D$, which vanishes throughout its particle region up to
$T+\epsilon<H$, gives a classical extension for the original data.
Section~\ref{cont:persistence} then proves its smoothness on every
compact slab; that conclusion was not imported from the cited theorem.

Finally, if $f_0=0$, the constraints make both initial fields
divergence-free.  Lemma~\ref{cont:vacuum} directly gives the global
solution with $f=0$, and classical uniqueness applies to it.
This proves the local existence, uniqueness, persistence, and
bounded-momentum continuation assertions of
Proposition~\ref{prop:continuation}.
